% manuscript-sha256: e68bf550e9a6da50d05af1fd9cf82ea2ef4b9550aa5eaba3bd0b987a798ecf74
% Options for packages loaded elsewhere
\PassOptionsToPackage{unicode}{hyperref}
\PassOptionsToPackage{hyphens}{url}
\documentclass[
]{article}
\usepackage{xcolor}
\usepackage[margin=1in]{geometry}
\usepackage{amsmath,amssymb}
\setcounter{secnumdepth}{-\maxdimen} % remove section numbering
\usepackage{iftex}
\ifPDFTeX
  \usepackage[T1]{fontenc}
  \usepackage[utf8]{inputenc}
  \usepackage{textcomp} % provide euro and other symbols
\else % if luatex or xetex
  \usepackage{unicode-math} % this also loads fontspec
  \defaultfontfeatures{Scale=MatchLowercase}
  \defaultfontfeatures[\rmfamily]{Ligatures=TeX,Scale=1}
\fi
\usepackage{lmodern}
\ifPDFTeX\else
  % xetex/luatex font selection
  \setmainfont[]{Times New Roman}
  \setmonofont[]{Consolas}
  \setmathfont[]{Latin Modern Math}
\fi
% Use upquote if available, for straight quotes in verbatim environments
\IfFileExists{upquote.sty}{\usepackage{upquote}}{}
\IfFileExists{microtype.sty}{% use microtype if available
  \usepackage[]{microtype}
  \UseMicrotypeSet[protrusion]{basicmath} % disable protrusion for tt fonts
}{}
\makeatletter
\@ifundefined{KOMAClassName}{% if non-KOMA class
  \IfFileExists{parskip.sty}{%
    \usepackage{parskip}
  }{% else
    \setlength{\parindent}{0pt}
    \setlength{\parskip}{6pt plus 2pt minus 1pt}}
}{% if KOMA class
  \KOMAoptions{parskip=half}}
\makeatother
\setlength{\emergencystretch}{3em} % prevent overfull lines
\providecommand{\tightlist}{%
  \setlength{\itemsep}{0pt}\setlength{\parskip}{0pt}}
\usepackage{bookmark}
\IfFileExists{xurl.sty}{\usepackage{xurl}}{} % add URL line breaks if available
\urlstyle{same}
\hypersetup{
  hidelinks,
  pdfcreator={LaTeX via pandoc}}

\author{}
\date{}

\allowdisplaybreaks
\begin{document}

\begin{center}
{\LARGE\bfseries Candidate selection with residual and constraint diagnostics\par}
Jeremy H. Carroll\par
2026-10-23\par
Version 2.0.4\par
\end{center}

\subsection{Abstract}\label{abstract}

We specify a finite candidate-selection rule using a quotient residual,
constraint violations, a closure residual, interface discrepancies and a
secondary alignment statistic. The quotient and closure residuals
satisfy an operator-norm inequality, so these channels are not
algebraically independent. Candidates must satisfy literal non-increase
of the quotient residual and an exact ratio ceiling before entering a
two-stage tolerance ranking. A separate inequality can screen
representation-expanding proposals. We establish the residual bound,
explain the ranking rule and give counterexamples to stronger
interpretations: a ratio ceiling does not preserve constraints against
deletion, and greedy candidate selection does not establish convergence.
The implementation provides exact rational decision routines alongside a
numerical interface.

\subsection{1. Introduction}\label{introduction}

A repair proposal can improve one diagnostic while worsening another. A
selection rule must therefore state which comparisons are mandatory and
which are permitted trade-offs. The rule studied here is a finite
filtering and ranking procedure. Its acceptance conditions are
definitions of a policy; they are not theorems that selected proposals
solve the underlying problem. The geometric setup uses a finite cochain
complex, while the optimization ingredients are finite-dimensional norm
minimization and ordered comparisons {[}1,2{]}.

\subsection{2. State variables and the residual
inequality}\label{state-variables-and-the-residual-inequality}

Let \(C^0\xrightarrow{B}C^1\xrightarrow{D}C^2\) be finite-dimensional
real coordinate spaces with \(DB=0\). For a field \(\omega\in C^1\),
define

\[
\Phi(\omega)=\min_a\|\omega-Ba\|_1,\qquad
\Gamma(\omega)=\|D\omega\|_1.
\]

Additional channels are a nonnegative clause-violation count \(S\), an
interface discrepancy

\[
J=\sum_{(i,j)}\|r_{ij}(\omega_i)-r_{ji}(\omega_j)\|_1,
\]

and a supplied alignment statistic \(X\in[0,1]\). If \(X\) is obtained
by a normalized residual construction, we use the separate convention
\(X=0\) when \(\Phi=0\). A tuple \((\Phi,S,\Gamma,J,X)\) does not by
itself establish that its entries are realized by a particular field or
constraint system.

\textbf{Proposition 2.1.} For the induced sum-norm operator bound
\(L=\|D\|_{1\to1}\),

\[
\Gamma(\omega)\le L\Phi(\omega).
\]

\textbf{Proof.} For every \(a\), \(D(\omega-Ba)=D\omega\), and hence
\(\|D\omega\|_1\le L\|\omega-Ba\|_1\). Minimization gives the result. In
finite dimensions \(\operatorname{im}B\) is closed and the residual
minimum is attained. In particular \(\Phi=0\) implies \(\Gamma=0\).

Different channels can nevertheless distinguish different defects. A
logical assignment can violate a clause even when the cochain is exact.
A directed cycle without faces has \(\Gamma=0\) for every field but can
have \(\Phi>0\). On a filled directed triangle,

\[
\Phi(\omega)=|\omega_1+\omega_2+\omega_3|=\Gamma(\omega).
\]

The lower bound is the triangle inequality for every representative; a
potential can move the whole cycle sum onto a single edge, attaining it.
These examples show distinct diagnostic roles, rather than complete
independence.

\subsection{3. Admissible candidates}\label{admissible-candidates}

For abstract nonnegative diagnostic pairs define

\[
R(\Phi,\Gamma)=
\begin{cases}
\Gamma/\Phi,&\Phi>0,\\
0,&\Phi=\Gamma=0,\\
+\infty,&\Phi=0<\Gamma.
\end{cases}
\]

The last case is incompatible with a realized field in the fixed complex
above. It is retained to reject inconsistent supplied data. Likewise
\(\Phi=0<X\) is rejected under the stated alignment convention.

Fix \(R_{\max}>0\). A candidate with primed diagnostics must satisfy

\[
\Phi'\le\Phi,\qquad R(\Phi',\Gamma')\le R_{\max}.
\]

These are literal comparisons. The ranking tolerance introduced below
does not relax them. For positive \(\Phi'\) the ratio comparison can be
evaluated exactly as \(\Gamma'\le R_{\max}\Phi'\) without division.

\textbf{Counterexample 3.1: deletion.} Removing one edge and the
incident face of a filled triangle leaves a path. Every field on that
path is exact and there are no face residuals. Both diagnostics become
zero and the ratio test passes. Consequently this ratio condition does
not detect constraint deletion. Any required preservation of the
carrier, clauses or interfaces is a separate condition on the proposed
transformation.

For a candidate that expands a representation, an optional additional
condition is

\[
(\Phi-\Phi')+\frac12(S-S')>c_{\mathrm{dim}},
\]

where \(c_{\mathrm{dim}}\ge0\) is a specified dimension penalty. This is
another policy condition, not a consequence of the residual inequality.

\subsection{4. Global tolerance ranking}\label{global-tolerance-ranking}

Let \(A\) be a finite, nonempty set of candidates that passed the
required checks. For candidate \(o\), let its nonnegative execution cost
be \(c_o\) and define

\[
F_o=\Phi_o+\tfrac12S_o+\tfrac15\Gamma_o+\tfrac1{10}J_o+c_o.
\]

The coefficients are specified choices. No mathematical optimality of
these particular weights is asserted. For a tolerance \(\epsilon>0\)
form

\[
F_{\min}=\min_{o\in A}F_o,\qquad
C=\{o\in A:F_o\le F_{\min}+\epsilon\},
\]

\[
X_{\min}=\min_{o\in C}X_o,\qquad
T=\{o\in C:X_o\le X_{\min}+\epsilon\}.
\]

Select the lexicographically smallest
\((F_o,X_o,\operatorname{name}(o))\) among candidates in \(T\), using a
fixed ordering of names. This is a global set-based rule. Pairwise
approximate equality is not used as a comparator. Identical records can
share a name; distinct selected transformations require distinct names
or another declared final ordering.

\textbf{Proposition 4.1.} The selected candidate has
\(F_o\le F_{\min}+\epsilon\) and \(X_o\le X_{\min}+\epsilon\). It need
not minimize either coordinate exactly.

\textbf{Proof.} Both bands contain a minimizer and are nonempty. The
final choice is made within \(T\subseteq C\), so both inequalities
follow from membership.

For example, with \(\epsilon=1/10\), scores and alignment values
\((0,1/2)\) and \((1/20,0)\) both enter the score band, but only the
second enters the alignment band. It is selected although its score is
larger. This is the intended tolerance rule. It does not permit the
primary residual to increase, because the admission comparison was
applied before ranking.

\subsection{5. Implementation and formal
verification}\label{implementation-and-formal-verification}

The exact rational methods of
\nolinkurl{code/mod_constraint_regulation_core.py} implement signature
validation, the cross-multiplied ratio comparison, the expansion
condition and the global ranking rule. They reject malformed rational
inputs. The numerical interface uses floating-point diagnostics and
rejects nonfinite states or scores. Numerical evaluations are not exact
proofs about the field from which those diagnostics are said to
originate.

\nolinkurl{lean4/Regulate.lean} verifies rational comparison statements,
selected ranking cases and classifications of rejected candidates. It
does not construct a general cochain complex or prove Proposition 2.1;
that proposition has the prose proof in Section 2. The bundled census
records the checked declaration set. Tests exercise the reference
implementation, including invalid inputs and the constraint-deletion
limitation.

\subsection{6. Limitations and
discussion}\label{limitations-and-discussion}

Selection uses a finite supplied candidate set. It does not establish
that this set contains a feasible repair. The policy does not certify
the diagnostic computation, preserve constraints against deletion or
guarantee termination. Even literal non-increase of a real residual
allows repeated equal-score steps and does not imply finite termination.
Applications needing an absolute closure bound must specify one
separately from the ratio ceiling.

\subsection{References}\label{references}

\begin{enumerate}
\def\labelenumi{\arabic{enumi}.}
\tightlist
\item
  S. Boyd and L. Vandenberghe, \emph{Convex Optimization}, Cambridge
  University Press,

  \begin{enumerate}
  \def\labelenumii{\arabic{enumii}.}
  \setcounter{enumii}{2003}
  \tightlist
  \item
    \href{https://web.stanford.edu/~boyd/cvxbook/}{Author-hosted book}.
  \end{enumerate}
\item
  A. Hatcher, \emph{Algebraic Topology}, Cambridge University Press,
  2002, Sections 2.1 and 3.1.
  \href{https://pi.math.cornell.edu/~hatcher/AT/ATpage.html}{Author-hosted
  book}.
\end{enumerate}

\end{document}
